\chapter{Vaginal Dryness}

\section*{Definition}
Insufficient vaginal lubrication causing discomfort, irritation, and pain during sexual activity or daily activities, most commonly due to decreased estrogen levels.

\section*{Pathophysiology}
Estrogen maintains vaginal mucosal thickness, elasticity, and lubrication via cervical and Bartholin gland secretions. Estrogen decline (menopause, breastfeeding, medications) leads to mucosal atrophy, decreased blood flow, reduced secretions, and increased vaginal pH predisposing to infection. Vaginal dryness is symptomatic of the genitourinary syndrome of menopause (GSM).

\section*{Etiology}
\begin{itemize}
  \item Menopause (natural or surgical)
  \item Breastfeeding or postpartum period
  \item Medication effects (antihistamines, antidepressants, chemotherapy, tamoxifen)
  \item Sjogren syndrome (autoimmune)
  \item Radiation therapy to pelvis
  \item Hormonal contraception (some types)
  \item Perimenopause
  \item Oophorectomy (surgical removal of ovaries)
  \item Anxiety or inadequate arousal
  \item Chemotherapy or endocrine therapy for breast cancer
\end{itemize}

\section*{Indications for Evaluation}
Seek care if:
\begin{itemize}
  \item Severe or worsening symptoms
  \item Vaginal dryness with painful intercourse, itching, burning, or recurrent urinary tract infections
  \item Accompanied by other concerning signs
\end{itemize}

\section*{Related Symptoms}
\begin{itemize}
  \item Painful sex (dyspareunia)
  \item Vaginal itching
  \item Urinary frequency
  \item Recurrent UTIs
  \item Vaginal irritation
\end{itemize}
\textbf{Note:} Always consider serious underlying conditions when evaluating persistent vaginal dryness.

\section*{Self-Care / Home Management}
\begin{itemize}
  \item Rest and avoid activities that may aggravate the condition.
  \item Maintain adequate hydration and a balanced diet.
  \item Over-the-counter remedies may provide symptomatic relief (consult a pharmacist or doctor).
  \item Monitor symptoms and seek professional advice if they persist or worsen.
  \item Avoid self-medicating with prescription medications without medical guidance.
\end{itemize}

\section*{Diagnostic Approach}
\begin{itemize}
  \item A thorough history and physical examination are the first steps in evaluation.
  \item Laboratory tests (blood work, urinalysis) may be ordered based on clinical suspicion.
  \item Imaging studies (X-ray, ultrasound, CT, MRI) may be indicated when structural pathology is suspected.
  \item Specialist referral is arranged when the diagnosis remains uncertain or requires specific expertise.
\end{itemize}

\section*{Treatment / Management Overview}
\begin{itemize}
  \item Treatment targets the underlying cause identified through diagnostic evaluation.
  \item Supportive care and symptom management are provided while awaiting definitive therapy.
  \item Medications, lifestyle modifications, and physical therapies may be prescribed as appropriate.
  \item Follow-up ensures treatment effectiveness and allows timely adjustments to the care plan.
\end{itemize}

\clearpage